\documentclass[reqno]{amsart}

\usepackage{amsmath,amssymb,amsthm,amsfonts}
\usepackage{mathrsfs}
\usepackage{hyperref}
\usepackage{enumitem}
\usepackage{mathtools}
\usepackage{booktabs}

\hypersetup{
  colorlinks=true,
  linkcolor=blue,
  citecolor=blue,
  urlcolor=blue
}

\numberwithin{equation}{section}

\theoremstyle{plain}
\newtheorem{theorem}{Theorem}[section]
\newtheorem{proposition}[theorem]{Proposition}
\newtheorem{corollary}[theorem]{Corollary}
\newtheorem{conjecture}[theorem]{Conjecture}
\newtheorem{lemma}[theorem]{Lemma}

\theoremstyle{definition}
\newtheorem{definition}[theorem]{Definition}
\newtheorem{axiom}{Axiom}

\theoremstyle{remark}
\newtheorem{remark}[theorem]{Remark}

%% Macros
\newcommand{\R}{\mathbb{R}}
\newcommand{\T}{\mathbb{T}}
\newcommand{\dm}{\mathfrak{D}}
\newcommand{\dmcat}{\mathbf{dm3}}
\newcommand{\Orb}{\Gamma}
\newcommand{\cN}{\mathcal{N}}
\newcommand{\cB}{\mathcal{B}}
\newcommand{\cS}{\mathcal{S}}
\newcommand{\cL}{\mathcal{L}}
\newcommand{\cR}{\mathcal{R}}
\newcommand{\cV}{\mathcal{V}}
\newcommand{\cM}{\mathcal{M}}
\newcommand{\cU}{\mathcal{U}}
\newcommand{\cP}{\mathcal{P}}
\newcommand{\cD}{\mathcal{D}}
\newcommand{\eps}{\varepsilon}
\newcommand{\dg}{d_g}
\newcommand{\norm}[1]{\left\|#1\right\|}
\newcommand{\abs}[1]{\left|#1\right|}
\DeclareMathOperator{\lcm}{lcm}
\DeclareMathOperator{\gcd}{gcd}
\DeclareMathOperator{\Fix}{Fix}
\DeclareMathOperator{\Hess}{Hess}
\DeclareMathOperator{\supp}{supp}

\begin{document}

\title[Generative Contact Mechanics]{%
  Generative Contact Mechanics:\\
  A Geometric Framework for Dissipative Systems\\
  with Structured Limit Cycles}

\author{Pablo Nogueira Grossi}
\address{Independent Researcher, Brazil}
\email{pablogrossi@hotmail.com}

\keywords{contact geometry, limit cycles, dissipative dynamical systems,
  operator algebra, structural stability, resonance, stochastic stability,
  nonequilibrium dynamics, Lyapunov functions}

\subjclass[2020]{37C10, 37C27, 37D30, 53D10, 60H10, 18A30}

\begin{abstract}
We introduce a geometric framework for dissipative dynamical systems
whose limit cycles carry structured phase, multi-scale coherence, and
stochastic stability. The central object is the \emph{dm3 system}---a
smooth Riemannian manifold equipped with a hyperbolic limit cycle, a
Lyapunov function, and a stochastic extension---formalized through eight
axioms. On this foundation we construct a complete operator algebra
(the $g$-, $L$-, $R$-, and $U$-operators) governing expansion, coherence,
resonance, and unification, and we embed the entire structure in contact
geometry.

Four main results are established. \textbf{Theorem~A}: any dm3 system
admits a unique exponentially stable hyperbolic limit cycle with a
topological invariant (the winding integral $\oint_{\Orb}\lambda$) and
stochastic stability below an embodiment threshold~$\tau$.
\textbf{Theorem~B}: the category $\dmcat$ is closed under a unification
operator~$\cU$, with the quantitative prediction $\tau_{12}\le\min(\tau_1,\tau_2)$:
unification weakens noise tolerance.
\textbf{Theorem~C}: every dm3 system near its limit cycle is locally
contact-diffeomorphic to a universal three-parameter normal form
$(\mu_{\max},\omega,\beta)$.
\textbf{Theorem~D}: the full operator algebra is $C^1$-structurally stable
with explicit stability radius
$\varepsilon_0 = \tfrac{1}{2}\abs{\mu_{\max}}\cdot\sup_{\Orb}\norm{\Hess V}^{-1}$.
Complete verification on an explicit toy model appears in the companion
paper~\cite{Paper2}.
\end{abstract}

\dedicatory{Dedicated to my children Vic~(R.I.P.), Giulia,
 Alice~(R.I.P.), Sarah~(R.I.P.), and David.\\[4pt]
\emph{Once tiny, always strong.}}

\maketitle
\tableofcontents

%%=============================================================
\section{Introduction}\label{sec:intro}
%%=============================================================

\subsection{The problem: dissipative systems with structured limit cycles}

Dissipative dynamical systems with hyperbolic limit cycles appear throughout
the natural sciences---in biology as oscillatory cellular and neural
processes, in physics as driven nonequilibrium systems, in engineering as
controlled oscillators, and in mathematics as paradigmatic examples of
nonlinear dynamics. The classical theory provides excellent tools for
analyzing a single limit cycle in isolation: Floquet theory characterizes
transverse stability, Lyapunov theory gives quantitative convergence rates,
and stochastic extensions handle noise.

What is less well understood is how limit cycles interact, unify, and
organize into multi-scale coherent structures while remaining robust under
perturbation. This paper develops a geometric framework---generative contact
mechanics---for precisely this class of behavior. The central insight is
that the correct geometric home for dissipative systems with structured limit
cycles is not symplectic geometry (which forbids attractors) but contact
geometry, where dissipation is encoded in the geometry itself.

\subsection{The contact geometric approach}

A hyperbolic limit cycle is an attractor: nearby trajectories converge to it
exponentially. This is incompatible with Hamiltonian mechanics, where
Liouville's theorem guarantees that the phase-space volume is preserved and
no attractors can exist on compact symplectic manifolds. Standard Lagrangian
mechanics is similarly insufficient: the Euler--Lagrange equations are
time-reversible and do not admit attractors without additional dissipation
terms added by hand.

Contact geometry provides a natural resolution. A contact manifold
$(M^{2n+1},\alpha)$ carries a Reeb vector field and a contact Hamiltonian
formalism in which dissipation is built into the geometric structure: the
contact Hamiltonian $H$ evolves along trajectories rather than remaining
constant, and the action variable $z\in\R$ encodes accumulated dissipation.
This makes contact geometry the natural setting for dissipative systems
with limit cycle attractors, stochastic stability thresholds, and
variational structure~\cite{Bravetti2017,deLeonLainz2019}.

The present paper formalizes this setting through the dm3 system and
operator algebra, and shows that the resulting framework is: (a)~geometrically
natural (contact normal form, Theorem~C), (b)~categorically coherent
(closure of $\dmcat$ under unification, Theorem~B), and (c)~physically
robust ($C^1$-structural stability, Theorem~D).

\subsection{Main results}

The four main theorems are stated informally here and proved in later
sections. Precise statements appear at the end of each relevant section.

\begin{theorem}[A: dm3 existence and stability]\label{thm:A}
Let $\dm=(S,g,Y,\Orb,V,\sigma)$ satisfy
Axioms~\ref{ax:orbit}--\ref{ax:hyperbolic}. Then:
\begin{enumerate}[label=\textup{(\roman*)}]
  \item $\Orb$ is a hyperbolic limit cycle of $Y$ with minimal period $T^*>0$.
  \item $\Orb$ is exponentially orbitally stable with rate $c=\abs{\mu_{\max}}$.
  \item For noise amplitude below the embodiment threshold $\tau$,
        $\mathbb{E}[V(X_t)]\to 0$.
  \item The winding integral $\oint_{\Orb}\lambda$ is a topological invariant
        preserved under all admissible deformations.
        \textup{(Proved in Proposition~\ref{prop:winding}.)}
\end{enumerate}
\end{theorem}

\begin{theorem}[B: closure under unification]\label{thm:B}
Let $\dm_1,\dm_2\in\dmcat$ be in $k$:$m$ resonance with admissible span.
Then:
\begin{enumerate}[label=\textup{(\roman*)}]
  \item The unification operator $\cU=U_3\circ U_2\circ U_1\circ R_3\circ R_2\circ R_1$
        is well-defined and $\cU(\dm_1,\dm_2)\in\dmcat$.
  \item The canonical invariants satisfy
        $T^*_{12}=kT^*_1/\gcd(k,m)$,\;
        $\mu_{\max,12}=\max(\mu_{\max,i})$,\;
        $\tau_{12}\le\min(\tau_1,\tau_2)$.
  \item \textbf{Unification weakens noise tolerance:}
        $\tau_{12}\le\min(\tau_1,\tau_2)$.
\end{enumerate}
\end{theorem}

\begin{theorem}[C: contact normal form]\label{thm:C}
Let $\dm$ satisfy Axioms~\ref{ax:orbit}--\ref{ax:hyperbolic}. In a tubular
neighborhood of $\Orb$ in $(M,\alpha)$, there exist contact coordinates
$(\rho,\theta,z)$ with $\Orb=\{\rho=0\}$ such that
\begin{align*}
  \dot\rho &= \mu_{\max}(1-e^{-\beta z})\rho + O(\rho^2),\\
  \dot\theta &= \omega + O(\rho),\\
  \dot z &= \omega - \abs{\mu_{\max}}\rho^2 e^{-\beta z} + O(\rho^3).
\end{align*}
Every dm3 system is locally contact-diffeomorphic to this normal form.
The parameters $(\mu_{\max},\omega,\beta)$ are the canonical invariant
triple.
\end{theorem}

\begin{theorem}[D: structural stability]\label{thm:D}
Let $\dm$ satisfy Axioms~\ref{ax:orbit}--\ref{ax:hyperbolic} with
$\mu_{\max}<0$. Define
\[
  \varepsilon_0 \coloneqq
  \frac{\abs{\mu_{\max}}}{2\bigl(1+\sup_{\Orb}\norm{\Hess V}_\infty\bigr)}.
\]
For any $C^1$ perturbation $F$ with $\norm{F}_{C^1}<\varepsilon_0$, and
any finite admissible operator sequence within parameter bounds, the output
$\dm'$ satisfies Axioms~\ref{ax:orbit}--\ref{ax:hyperbolic} with: a unique
hyperbolic limit cycle $\Orb'$ near $\Orb$; the same winding number;
canonical invariants varying by $O(\varepsilon_0)$; and boundaries $B_i'$
varying $C^1$-smoothly.
\end{theorem}

\subsection{Relation to existing work}\label{sec:related}

\emph{Contact Hamiltonian mechanics.}
Bravetti~\cite{Bravetti2017} and de~Le\'on--Lainz~\cite{deLeonLainz2019}
developed the modern formalism of contact Hamiltonian mechanics, showing
that contact geometry provides a variational framework for dissipative
systems. Gaset et al.~\cite{Gaset2020} extended this to Lagrangian
contact mechanics. The present paper differs from these works in three
respects: (a)~we focus on limit cycle attractors rather than fixed-point
attractors; (b)~we construct a categorical operator algebra acting on
contact systems; and (c)~we prove structural stability of the operator
algebra with an explicit stability radius.

\emph{Metriplectic systems.}
Morrison~\cite{Morrison1986} introduced the metriplectic framework for
systems combining Hamiltonian and dissipative dynamics. The dm3 contact
formulation is related but distinct: metriplectic structure separates the
Poisson and metric brackets, while our contact structure encodes dissipation
in the contact form itself through the action variable $z$.

\emph{Normally hyperbolic invariant manifolds.}
Hirsch, Pugh, and Shub~\cite{HPS} established the persistence and
$C^r$-regularity of normally hyperbolic invariant manifolds under
perturbation. Theorem~D uses this theory as a foundational input.
The novelty is combining NHIM persistence with the operator algebra
and establishing the quantitative stability radius $\varepsilon_0$.

\emph{Arnold tongues and coupled oscillators.}
The resonance structure of the $R$-operators (period resonance, Arnold
tongues, mode locking) is classical~\cite{Arnold,Pikovsky}. The contribution
here is embedding these phenomena in a categorical structure: resonance is
the precondition for a well-defined pushout in $\dmcat$. The prediction
$\tau_{12}\le\min(\tau_i)$ is new.

\subsection{Organization}

Section~\ref{sec:background} fixes notation and background.
Section~\ref{sec:dm3} defines dm3 systems and proves Theorem~A.
Section~\ref{sec:operators} constructs the operator algebra and proves
Theorem~B.
Section~\ref{sec:boundary} establishes boundary operators and the
operator grammar.
Section~\ref{sec:contact} develops the contact physics and proves
Theorem~C.
Section~\ref{sec:stability} proves Theorem~D.
Section~\ref{sec:discussion} discusses open problems and connections.
Appendices give the categorical and lemma details.

%%=============================================================
\section{Background}\label{sec:background}
%%=============================================================

\subsection{Riemannian manifolds and flows}

Throughout, $(S,g)$ denotes a smooth, connected, complete Riemannian
manifold of dimension $n\ge 2$. We write $T_xS$ for the tangent space
at $x$, $\mathfrak{X}(S)$ for smooth vector fields, and $\Omega^k(S)$
for smooth $k$-forms. For $X\in\mathfrak{X}(S)$ complete, the flow
$\phi^t\colon S\to S$ satisfies $\phi^0=\mathrm{id}$ and
$\tfrac{d}{dt}\phi^t(x)=X(\phi^t(x))$.

\subsection{Hyperbolic limit cycles and Floquet theory}

\begin{definition}\label{def:limitcycle}
A \emph{limit cycle} of $Y\in\mathfrak{X}(S)$ is a compact connected
invariant submanifold $\Orb\cong S^1$ with minimal period $T^*>0$.
It is \emph{hyperbolic} if all Floquet multipliers except the trivial
multiplier $\mu_0=1$ satisfy $\abs{\mu_j}<1$ for $j=1,\ldots,n-1$.
The \emph{maximal transverse Lyapunov exponent} is
$\mu_{\max}\coloneqq\max_j\tfrac{1}{T^*}\log\abs{\mu_j}<0$.
\end{definition}

The stable manifold theorem~\cite{HPS} applied to $\Orb$ gives: if
$\Orb$ is hyperbolic, then $\partial\cB(\Orb)=\mathcal{W}^s(\Lambda)$
for some invariant set $\Lambda\subset S\setminus\cB(\Orb)$.

\subsection{Lyapunov functions}

\begin{definition}\label{def:lyapunov}
A smooth function $V\colon S\to\R_{\ge 0}$ is a \emph{Lyapunov function}
for $\Orb$ if $V(x)=0\iff x\in\Orb$ and
$\dot V(x)\coloneqq g(\nabla V(x),Y(x))<0$ for $x\notin\Orb$.
The standard choice is $V=\dg(\cdot,\Orb)^2$.
\end{definition}

\subsection{Stochastic differential equations}

For the It\^o SDE $dX_t=Y(X_t)\,dt+\sigma(X_t)\,dW_t$ on $S$,
the generator is
\[
  \cL f(x) = g(\nabla f(x),Y(x))
             + \tfrac12\mathrm{tr}\bigl(\sigma(x)^\top\Hess_f(x)\,\sigma(x)\bigr).
\]
Standard stochastic Lyapunov theory~\cite{Hasminski} gives: if
$\cL V\le -cV+\kappa\norm\sigma^2$ for $c,\kappa>0$, then
$\mathbb{E}[V(X_t)]\to 0$ for noise amplitude below $\tau=\sqrt{c/\kappa}$.

\subsection{Contact manifolds}

\begin{definition}\label{def:contact}
A \emph{contact manifold} $(M^{2n+1},\alpha)$ is a smooth manifold
with a 1-form $\alpha$ satisfying $\alpha\wedge(d\alpha)^n\ne 0$
everywhere. The \emph{Reeb vector field} $R_\alpha$ is the unique
vector field with $\iota_{R_\alpha}d\alpha=0$ and $\alpha(R_\alpha)=1$.
For a smooth $H\colon M\to\R$, the \emph{contact Hamiltonian vector
field} $X_H$ is the unique vector field satisfying
$\iota_{X_H}d\alpha=dH-(R_\alpha H)\alpha$ and $\alpha(X_H)=-H$.
\end{definition}

See Geiges~\cite{Geiges} and Bravetti~\cite{Bravetti2017} for background.

%%=============================================================
\section{The dm3 operator}\label{sec:dm3}
%%=============================================================

\subsection{Definition}

\begin{definition}[dm3 system]\label{def:dm3}
A \emph{dm3 system} is a tuple $\dm=(S,g,Y,\Orb,V,\sigma)$ where
$(S,g)$ is a smooth complete Riemannian manifold, $Y\in\mathfrak{X}(S)$
is a smooth vector field with complete flow $\phi^t$, $\Orb\subset S$
is a hyperbolic limit cycle, $V\colon S\to\R_{\ge 0}$ is a smooth
Lyapunov function for $\Orb$, and $\sigma\colon S\to\R^{n\times m}$
is a smooth noise coefficient. The \emph{dm3 operator} is
$\mathcal{A}_{\mathrm{dm3}}\coloneqq\phi^T$ for $T=T^*/4$.
\end{definition}

\subsection{The eight axioms}\label{subsec:axioms}

\begin{axiom}[Orbit identity]\label{ax:orbit}
There exist four distinct points $p_I,p_{VI},p_{III},p_{VII}\in\Orb$
traversed cyclically:
$\phi^{T^*/4}(p_I)=p_{VI}$, $\phi^{T^*/4}(p_{VI})=p_{III}$,
$\phi^{T^*/4}(p_{III})=p_{VII}$, $\phi^{T^*/4}(p_{VII})=p_I$.
\end{axiom}

\begin{axiom}[Generative closure]\label{ax:closure}
$\Orb$ is compact, connected, and invariant:
$\phi^t(\Orb)=\Orb$ for all $t\in\R$.
\end{axiom}

\begin{axiom}[Structural primacy]\label{ax:primacy}
There exists a neighborhood $\cN(\Orb)$ such that
$\dot V(x)=g(\nabla V(x),Y(x))<0$ for all $x\in\cN(\Orb)\setminus\Orb$.
\end{axiom}

\begin{axiom}[Perturbation response]\label{ax:response}
There exist $c>0$ and $\cN(\Orb)$ such that
$\dot V(x)\le -c\,V(x)$ for all $x\in\cN(\Orb)$.
\end{axiom}

\begin{axiom}[Noise-as-fuel]\label{ax:noise}
The SDE generator satisfies
$\cL V(x)\le -c\,V(x)+\kappa\norm{\sigma(x)}^2$
for constants $c,\kappa>0$.
\end{axiom}

\begin{axiom}[Non-collapse]\label{ax:noncollapse}
$Y(x)\ne 0$ for all $x\in\Orb$, and $\phi^t$ has minimal period
$T^*>0$ on $\Orb$.
\end{axiom}

\begin{axiom}[Non-coercion]\label{ax:smooth}
$Y$ is $C^\infty$ on $S$. No discontinuous projection, clamping,
or hard boundary enforcement appears in the definition of $Y$.
\end{axiom}

\begin{axiom}[Hyperbolicity]\label{ax:hyperbolic}
All Floquet multipliers of $\Orb$ except the trivial one satisfy
$\abs{\mu_j}<1$ for $j=1,\ldots,n-1$.
\end{axiom}

\begin{remark}
These eight axioms are logically independent. The original ten axioms
are reduced here by: merging the stochastic pair into
Axiom~\ref{ax:noise}; and moving invariance under the $g$- and
$L$-operators and multi-scale compatibility to their respective
sections as theorems.
\end{remark}

\subsection{Orbital stability}

\begin{proposition}[Exponential orbital stability]\label{prop:orbital}
Let $\dm$ satisfy Axioms~\ref{ax:orbit}--\ref{ax:hyperbolic}. Then
there exist $c>0$ and $\cN(\Orb)$ such that
$\dot V(x)\le -c\,V(x)$ for all $x\in\cN(\Orb)$.
\end{proposition}

\begin{proof}
Axiom~\ref{ax:hyperbolic} gives $\mu_{\max}<0$. By Floquet theory
(see~\cite{GH,Shub}), the linearized transverse flow contracts with
rate $\abs{\mu_{\max}}$. With $V=\dg(\cdot,\Orb)^2$, we get
$\dot V\le 2\mu_{\max}V+O(V^{3/2})$ in a tubular neighborhood.
Taking $\cN(\Orb)$ small enough absorbs the higher-order term,
giving $\dot V\le -c\,V$ with $c=\abs{\mu_{\max}}$.
\end{proof}

\subsection{Embodiment threshold and stochastic stability}

\begin{proposition}[Stochastic Lyapunov condition]\label{prop:stoch}
If Axiom~\ref{ax:noise} holds with constants $c,\kappa$, then for
noise amplitude $\norm\sigma<\tau\coloneqq\sqrt{c/\kappa}$:
\[
  \mathbb{E}[V(X_t)]\le e^{-ct}V(X_0)+\frac{\kappa}{c}\norm\sigma^2.
\]
\end{proposition}

\begin{proof}
Standard; see Has'mi\u{n}ski\u{\i}~\cite{Hasminski}, Chapter~4.
\end{proof}

\begin{definition}[Embodiment threshold]\label{def:threshold}
The \emph{embodiment threshold} is
$\tau\coloneqq\sup\{\sigma_0:\cL V(x)<0\text{ whenever }
\norm{\sigma(x)}<\sigma_0\}=\sqrt{c/\kappa}$.
Below $\tau$: noise reinforces the orbit. Above $\tau$: noise
disrupts it.
\end{definition}

\begin{definition}[Canonical invariant triple]\label{def:triple}
$\iota(\dm)\coloneqq(T^*,\mu_{\max},\tau)\in
\R_{>0}\times\R_{<0}\times\R_{>0}$.
\end{definition}

\subsection{Proof of Theorem~A}

\begin{proof}[Proof of Theorem~\ref{thm:A}]
Claims (i)--(ii): Axioms~\ref{ax:orbit}, \ref{ax:closure},
\ref{ax:noncollapse}, \ref{ax:hyperbolic} give a compact connected
hyperbolic limit cycle with $T^*>0$. Proposition~\ref{prop:orbital}
gives exponential stability with $c=\abs{\mu_{\max}}$.
Claim~(iii): Proposition~\ref{prop:stoch} with $\tau=\sqrt{c/\kappa}$.
Claim~(iv): Proposition~\ref{prop:winding} in
Section~\ref{sec:contact}.
\end{proof}

%%=============================================================
\section{The operator algebra}\label{sec:operators}
%%=============================================================

Throughout, $\dm=(S,g,Y,\Orb,V,\sigma)$ is a fixed dm3 system
satisfying Axioms~\ref{ax:orbit}--\ref{ax:hyperbolic}.

\subsection{\texorpdfstring{$g$}{g}-series operators: expansion}
\label{subsec:g}

\subsubsection{\texorpdfstring{$g_1$}{g1}: local basin expansion}

\begin{definition}[$g_1$]\label{def:g1}
Let $\rho\colon S\to[0,1]$ be a smooth cutoff with $\rho=0$ on
$\cN_\delta(\Orb)$ and $\rho=1$ outside $\cN_{2\delta}(\Orb)$,
and let $h\colon\R_{\ge 0}\to\R_{\ge 0}$ be smooth with
$\supp h\subset(\delta,R)$ for $R>2\delta$. Define
\[
  Z(x)\coloneqq -\rho(x)\,h(V(x))\,\nabla V(x),\quad
  g_1\coloneqq\exp(Z)=\psi^1,
\]
where $\psi^t$ is the flow of $Z$.
\end{definition}

\begin{proposition}[$\Orb$-invariance of $g_1$]\label{prop:g1inv}
$g_1(\Orb)=\Orb$.
\end{proposition}
\begin{proof}
For $x\in\Orb$, $V(x)=0$, so $\rho(x)=0$ and $Z(x)=0$.
\end{proof}

\begin{proposition}[Basin expansion]\label{prop:g1expand}
The pushforward $Y^{(1)}\coloneqq(g_1)_*Y$ has Lyapunov function
$V^{(1)}=V\circ g_1^{-1}$ satisfying $\dot V^{(1)}\le -c^{(1)}V^{(1)}$
on $\cN^{(1)}(\Orb)\supset\cN(\Orb)$ with $c^{(1)}\ge c$.
The expanded system $\dm^{(1)}=(S,g,Y^{(1)},\Orb,V^{(1)},\sigma)$
is a dm3 system with $\tau^{(1)}\ge\tau$.
\end{proposition}

\begin{conjecture}[Global basin expansion]\label{conj:global}
For appropriate $h$ and $\delta$, the basin of attraction of $\Orb$
under $Y^{(1)}$ is strictly larger than under $Y$.
\end{conjecture}

\subsubsection{\texorpdfstring{$g_2$}{g2}: recursive multi-scale expansion}

\begin{definition}[Phase function and second-scale space]\label{def:phase}
The \emph{phase function} $\varphi\colon\cN(\Orb)\to\R/\mathbb{Z}$
assigns arc-length normalized phase, with $\varphi(p_I)=0$,
$\varphi(p_{VI})=1/4$, $\varphi(p_{III})=1/2$,
$\varphi(p_{VII})=3/4$. Set $S^{(2)}=\R/\mathbb{Z}$ with
projection $\Psi(x)=\varphi(x)$.
\end{definition}

\begin{definition}[$g_2$ and semi-conjugacy]\label{def:g2}
Let $\omega_2>0$, $f\colon S^{(2)}\to\R$ smooth. Define
$Y^{(2)}(\Phi)=\omega_2+f(\Phi)$ with flow $\psi^{(2),t}$, and
$g_2=\psi^{(2),1/\omega_2}$. The orbit $\Orb$ is
\emph{recursively invariant} if $\exists\,\lambda>0$ such that
\[
  \Psi(\phi^t(x)) = \psi^{(2),t/\lambda}(\Psi(x))
  \quad\forall\,x\in\Orb.
\]
Setting $\omega_2=1/T^*$ and $f\equiv 0$ achieves this with $\lambda=1$.
\end{definition}

\begin{conjecture}[Invariant torus stability]\label{conj:torus}
If the semi-conjugacy holds and $Y^{(2)}$ has a hyperbolic limit cycle,
the product system on $S\times S^{(2)}$ admits a stable invariant torus
(KAM-stable in the non-resonant case; mode-locked in the resonant case).
Proved in the toy model in~\cite{Paper2}.
\end{conjecture}

\subsubsection{\texorpdfstring{$g_3$}{g3}: cross-domain expansion and the category \texorpdfstring{$\dmcat$}{dm3}}

\begin{definition}[Category $\dmcat$]\label{def:cat}
The \emph{category} $\dmcat$ has objects: dm3 systems satisfying
Axioms~\ref{ax:orbit}--\ref{ax:hyperbolic}; morphisms: smooth maps
$f\colon S_1\to S_2$ satisfying:
\begin{enumerate}[label=\textup{(M\arabic*)},leftmargin=2em]
  \item $f(\Orb_1)=\Orb_2$,
  \item $f\circ\phi_1^t=\phi_2^{\lambda t}\circ f$ for some $\lambda>0$,
  \item $\alpha\,V_1(x)\le V_2(f(x))\le\beta\,V_1(x)$ for
        $0<\alpha\le\beta<\infty$,
  \item $\norm{\sigma_2(f(x))}\le C\norm{\sigma_1(x)}$ and
        $\tau_2\ge\alpha\tau_1/C$;
\end{enumerate}
composition: function composition; identity: $\mathrm{id}_S$.
\end{definition}

\begin{proposition}[Morphisms preserve dm3 structure]\label{prop:morphism}
If $f\colon\dm_1\to\dm_2$ is a dm3 morphism, then $\dm_2\in\dmcat$.
\end{proposition}
\begin{proof}
Axiom~\ref{ax:orbit}: (M1) and (M2) map phase points to phase points.
Axiom~\ref{ax:closure}: $f(\Orb_1)=\Orb_2$ is invariant by (M2).
Axioms~\ref{ax:primacy}--\ref{ax:response}: (M3) gives a Lyapunov
function for $\Orb_2$ with rate $c'\ge\alpha c/\beta$.
Axiom~\ref{ax:noise}: (M4). Axioms~\ref{ax:noncollapse}--\ref{ax:hyperbolic}:
smoothness of $f$ and (M1)--(M2) with $C^1$-diffeomorphism property
near $\Orb_1$.
\end{proof}

\begin{corollary}[Falsifiability of cross-domain claims]\label{cor:false}
The claim that the dm3 framework applies to a domain $D$ with state
space $S_D$ is equivalent to the existence of a morphism
$f\colon\dm_{\mathrm{base}}\to\dm_D$. This is a falsifiable condition.
\end{corollary}

\subsection{\texorpdfstring{$L$}{L}-operators: coherence and anti-collapse}
\label{subsec:L}

\subsubsection{\texorpdfstring{$L_1$}{L1}: local coherence}

\begin{definition}[$L_1$]\label{def:L1}
With two orbits $\Orb_1,\Orb_2$, inter-orbit distance
$D_{12}(x)=\dg(x,\Orb_1)+\dg(x,\Orb_2)$, and smooth cutoff
$\chi$ supported where $\min\{\dg(x,\Orb_i)\}<\eps$, define
$Z(x)=-\chi(x)\nabla D_{12}(x)$ and $L_1=\exp(Z)$.
\end{definition}

\begin{proposition}[$L_1$ orbit invariance]\label{prop:L1}
$L_1(\Orb_i)=\Orb_i$ for $i=1,2$. There exists $\eps'>\eps$ such
that $\dg(L_1(x),\Orb_i)\ge\dg(x,\Orb_i)$ for
$\dg(x,\Orb_i)\in[\eps,\eps']$.
\end{proposition}
\begin{proof}
$\chi=0$ on each $\Orb_i$, so $Z=0$ there. The gradient push
increases separation in the buffer zone by compactness.
\end{proof}

\begin{conjecture}[$L_1$ hyperbolicity]\label{conj:L1}
For sufficiently small $\norm{Z}_{C^1}$, both $\Orb_i$ remain
hyperbolic under $(L_1)_*Y_i$.
\end{conjecture}

\subsubsection{\texorpdfstring{$L_2$}{L2}: global coherence}

\begin{definition}[$L_2$]\label{def:L2}
The \emph{coherence defect}
$\mathcal{E}(x)=d_{S^{(2)}}(\Psi(\phi^{\Delta t}(x)),
\psi^{(2),\Delta t/\lambda}(\Psi(x)))^2$ measures failure of
semi-conjugacy. Set $\cV(x)=V(x)+\alpha_0\,\mathcal{E}(x)$,
$\alpha_0>0$, and $L_2=\exp(-\nabla\cV)$.
\end{definition}

\begin{proposition}[$L_2$ fixed points]\label{prop:L2}
$L_2(x)=x\iff x\in\Orb$. Moreover $\mathcal{E}|_{\Orb}=0$ and
$\nabla\cV|_{\Orb}=0$.
\end{proposition}
\begin{proof}
$\nabla\cV=0$ iff $\nabla V=0$ (i.e.\ $x\in\Orb$) and
$\nabla\mathcal{E}=0$. On $\Orb$ the semi-conjugacy holds
(Proposition~\ref{def:g2}) so $\mathcal{E}|_{\Orb}=0$ and
$\nabla\mathcal{E}|_{\Orb}=0$.
\end{proof}

\subsubsection{\texorpdfstring{$L_3$}{L3}: anti-collapse}

\begin{definition}[$L_3$]\label{def:L3}
For $\eps_0>0$, define the regularized potential
$U_{12}^\eps(x)=((\dg(x,\Orb_1)^2+\eps_0)^{-1}+(\dg(x,\Orb_2)^2+\eps_0)^{-1})$,
cutoff $\chi_{12}$ in the inter-orbit region,
$Q(x)=\chi_{12}(x)\nabla U_{12}^\eps(x)$,
and $L_3=\exp(\eta Q)$ for small $\eta>0$.
\end{definition}

\begin{proposition}[$L_3$ non-collapse and hyperbolicity]\label{prop:L3}
For sufficiently small $\eta$:
$\dg(L_3(\Orb_1),L_3(\Orb_2))\ge\dg(\Orb_1,\Orb_2)$,
and each $\Orb_i$ remains a hyperbolic limit cycle of $(L_3)_*Y_i$.
\end{proposition}
\begin{proof}
$Q$ points in the direction increasing $\dg(\Orb_1,\Orb_2)$ to first
order. Hyperbolicity is an open $C^1$ condition; $L_3$ is $C^1$-close
to the identity for small $\eta$.
\end{proof}

\subsection{\texorpdfstring{$R$}{R}-operators: resonance}
\label{subsec:R}

\subsubsection{\texorpdfstring{$R_1$}{R1}: detection}

\begin{definition}[$R_1$]\label{def:R1}
Systems $\dm_1,\dm_2$ are in \emph{$k$:$m$ period resonance}
if $kT_1^*=mT_2^*$; in \emph{Lyapunov resonance} if
$\mu_{\max,1}=\mu_{\max,2}$; in \emph{threshold resonance} if
$\tau_1=\tau_2$. Define
\[
  R_1(\dm_1,\dm_2)\coloneqq
  \{(k,m):kT_1^*=mT_2^*\}\cup\mathbf{1}[\mu_{\max,1}=\mu_{\max,2}]
  \cup\mathbf{1}[\tau_1=\tau_2].
\]
$R_1=\emptyset$ iff no resonance; only trivial coproduct
unification is then possible.
\end{definition}

\begin{proposition}[Genericity]\label{prop:R1}
Each resonance locus is codimension-$1$ in $\cP^2$ and has measure
zero. Within each Arnold tongue, $k$:$m$ resonance is structurally
stable under coupling.
\end{proposition}

\subsubsection{\texorpdfstring{$R_2$}{R2}: amplification}

\begin{definition}[$R_2$]\label{def:R2}
Phase mismatch:
$W_{k:m}(\theta_1,\theta_2)=(k\varphi_1(\theta_1)-m\varphi_2(\theta_2))^2$.
Define $Z_{k:m}=-\nabla_{(x_1,x_2)}W_{k:m}$ and
$R_2=\exp(Z_{k:m})$.
\end{definition}

\begin{proposition}[Phase locking and descent]\label{prop:R2}
$\Fix(R_2)\cap(\Orb_1\times\Orb_2)=\{k\varphi_1=m\varphi_2\}$,
consisting of $\lcm(k,m)$ points generically.
Along the flow: $\tfrac{d}{dt}W_{k:m}=-\norm{\nabla W_{k:m}}^2\le 0$.
\end{proposition}
\begin{proof}
Fixed points are zeros of $Z_{k:m}$, i.e.\ critical points of
$W_{k:m}$. On the torus $(\R/\mathbb{Z})^2$, the equation
$k\varphi_1-m\varphi_2=0$ has $\lcm(k,m)$ solutions generically.
Descent: $\dot W_{k:m}=-\norm{\nabla W_{k:m}}^2$ by definition.
\end{proof}

\subsubsection{\texorpdfstring{$R_3$}{R3}: stabilization}

\begin{definition}[$R_3$]\label{def:R3}
Resonance manifold:
$\cM_{k:m}=\{(x_1,x_2)\in\cN(\Orb_1)\times\cN(\Orb_2):
k\varphi_1(x_1)=m\varphi_2(x_2)\}$.
Define $R_3=\exp(-\nabla\dg(\cdot,\cM_{k:m})^2)$.
\end{definition}

\begin{proposition}[Exponential stability]\label{prop:R3}
With $U=d(\cdot,\cM_{k:m})^2$: $\dot U=-4U$ along $R_3$-flow,
so $U(t)=U(0)e^{-4t}$.
Both $\Orb_i$ are fixed by $R_3$.
\end{proposition}
\begin{proof}
$\dot U=-\norm{\nabla U}^2=-4d^2=−4U$ since $\norm{\nabla d}=1$
away from the cut locus. Orbit invariance: $\nabla d(\cdot,\cM_{k:m})$
is tangent to each $\Orb_i$ on $\Orb_i$, so the gradient flow moves
$\Orb_i$ along itself.
\end{proof}

\subsection{\texorpdfstring{$U$}{U}-operators: unification}
\label{subsec:U}

\subsubsection{Admissible spans and pushout existence}

\begin{definition}[Resonant orbit and admissible span]\label{def:span}
$\Orb_\lambda\coloneqq\cM_{k:m}\cap(\Orb_1\times\Orb_2)$
(consists of $\lcm(k,m)$ points). A span
$\dm_1\xleftarrow{p_1}\dm_\lambda\xrightarrow{p_2}\dm_2$
is \emph{admissible} if: (1)~$p_1,p_2$ are dm3 morphisms;
(2)~smooth embeddings; (3)~$\dm_\lambda$ has tubular neighborhoods
in each $S_i$; (4)~$V_1\circ p_1=V_2\circ p_2$ on $\dm_\lambda$.
\end{definition}

\begin{proposition}[Projections are morphisms]\label{prop:proj}
The projections $p_i\colon\cM_{k:m}\to S_i$ satisfy (M1)--(M4)
with $p_i(\Orb_\lambda)=\Orb_i$.
\end{proposition}
\begin{proof}
(M1): $p_i(\Orb_\lambda)=\Orb_i$ by definition.
(M2): $k$:$m$ resonance gives semi-conjugacy with $\lambda=k/m$.
(M3): $V_i\circ p_i$ equivalent near $\Orb_\lambda$ by transversality
of $\cM_{k:m}$ to $V_i$-level sets.
(M4): bounded noise coefficients.
\end{proof}

\begin{proposition}[Pushout existence]\label{prop:pushout}
If the span is admissible, the pushout
$\dm_{12}=\dm_1\sqcup_{\dm_\lambda}\dm_2$ exists in $\dmcat$,
unique up to dm3-isomorphism.
\end{proposition}
\begin{proof}
Conditions (1)--(3) give smooth gluing of $S_1$,$S_2$ along
$p_i(\cM_{k:m})$ by the smooth gluing theorem~\cite{Lee};
(4) gives smooth $V_{12}$. Hyperbolicity of $\Orb_{12}$ follows
from hyperbolicity of $\Orb_1,\Orb_2$ and transversality of the
gluing; uniqueness is categorical.
\end{proof}

\subsubsection{\texorpdfstring{$U_1$, $U_2$, $U_3$}{U1, U2, U3}}

\begin{definition}[$U_1$, $U_2$, $U_3$]\label{def:U}
$U_1(\dm_1,\dm_2)\coloneqq\dm_1\sqcup_{\dm_\lambda}\dm_2$.

The \emph{resonance correspondence}
$\cC_{12}=\{(x_1,x_2)\in S_1\times S_2:(x_1,x_2)\in\cM_{k:m}\}$
defines $U_2(x_1)=\pi_2(\pi_1^{-1}(x_1)\cap\cC_{12})$.

The \emph{synthesized orbit}
$\Orb_{\mathrm{syn}}=\omega\text{-limit set of generic trajectories}$;
$U_3=\exp(-\nabla\dg(\cdot,\Orb_{\mathrm{syn}})^2)$.
\end{definition}

\begin{proposition}[dm3 closure under $U_3$]\label{prop:closure}
After applying $U_3$, the system
$\dm_{12}=(S_{12},g_{12},Y_{12},\Orb_{\mathrm{syn}},V_{12},\sigma_{12})$
satisfies Axioms~\ref{ax:orbit}--\ref{ax:hyperbolic}.
\end{proposition}
\begin{proof}
$\Orb_{\mathrm{syn}}$ is a limit cycle by the $k$:$m$ resonance
collapsing the product torus; hyperbolicity from normal hyperbolicity
of $\cM_{k:m}$ (Proposition~\ref{prop:R3}). Axioms~\ref{ax:orbit}--\ref{ax:response}
hold by construction of $V_{12}$; Axiom~\ref{ax:noise} inherited
with $\tau_{12}\le\min(\tau_i)$; Axioms~\ref{ax:noncollapse}--\ref{ax:hyperbolic}
from the above.
\end{proof}

\subsection{Proof of Theorem~B}

\begin{proof}[Proof of Theorem~\ref{thm:B}]
(i) Well-definedness: $R_1\ne\emptyset$ (resonance assumed);
$R_2$ drives to phase locking (Proposition~\ref{prop:R2});
$R_3$ stabilizes $\cM_{k:m}$ (Proposition~\ref{prop:R3});
admissibility gives the pushout (Proposition~\ref{prop:pushout});
$U_1$--$U_3$ proceed as in Definition~\ref{def:U};
$\dm_{12}\in\dmcat$ by Proposition~\ref{prop:closure}.
(ii)--(iii) Period: $T_{12}^*=kT_1^*/\gcd(k,m)$ (smallest $T>0$
with $T/T_i^*\in\mathbb{Z}$). Stability: slowest contraction dominates.
Threshold: $c_{12}=\min(c_i)$, $\kappa_{12}=\max(\kappa_i)$,
giving $\tau_{12}=\sqrt{c_{12}/\kappa_{12}}\le\min(\tau_i)$.
\end{proof}

%%=============================================================
\section{Boundary operators and constraint geometry}\label{sec:boundary}
%%=============================================================

\subsection{\texorpdfstring{$B_1$}{B1}: identity boundary}

\begin{definition}[$B_1$]\label{def:B1}
$B_1\coloneqq\mathcal{W}^s(\Lambda)=\partial\cB(\Orb)$,
the stable manifold of the boundary invariant set $\Lambda$.
The critical Lyapunov level
$c^*\coloneqq\sup\{c:\{V<c\}\subset\cB(\Orb)\}$
gives inner approximation $\{V=c^*\}\subseteq B_1$.
\end{definition}

\begin{proposition}[$B_1$ invariance]\label{prop:B1}
$\phi^t(B_1)=B_1$ for all $t\in\R$.
\end{proposition}
\begin{proof}
$B_1=\mathcal{W}^s(\Lambda)$ is invariant by definition of stable manifolds.
\end{proof}

\subsection{\texorpdfstring{$B_2$}{B2}: transition boundary}

\begin{definition}[$B_2$ and Arnold tongues]\label{def:B2}
Parameter space: $\cP=\R_{>0}\times\R_{<0}\times\R_{>0}$.
Transition boundary in $\cP^2$:
\[
  B_2=\bigcup_{k,m}\{kT_1^*=mT_2^*\}
  \cup\{\mu_{\max,1}=\mu_{\max,2}\}\cup\{\tau_1=\tau_2\}.
\]
Arnold tongue:
$\mathcal{A}_{k:m}=\{\abs{kT_1^*-mT_2^*}<\eps_{k:m}(A)\}$.
Inside $\mathcal{A}_{k:m}$: $R$-operators valid. Outside: invalid.
\end{definition}

\subsection{\texorpdfstring{$B_3$}{B3}: collapse boundary}

\begin{definition}[$B_3$]\label{def:B3}
$U_{12}^\eps(x)=(\dg(x,\Orb_1)^2+\eps_0)^{-1}
+(\dg(x,\Orb_2)^2+\eps_0)^{-1}$. Define
$C_{12}^*=\sup\{c:\dot U_{12}^\eps>0\text{ whenever }U_{12}^\eps>c\}$
and $B_3=\{U_{12}^\eps=C_{12}^*\}$.
\end{definition}

\begin{proposition}[Smoothness of $B_3$]\label{prop:B3}
For generic $\eps_0$, $B_3$ is a smooth codimension-$1$ submanifold.
\end{proposition}
\begin{proof}
$U_{12}^\eps$ is smooth; Sard's theorem gives generic regular values;
the regular level set theorem applies.
\end{proof}

\subsection{Admissible regions and operator grammar}

\begin{definition}[Admissible regions]\label{def:admissible}
$\cD_S\coloneqq S\setminus(B_1\cup B_3)$;\quad
$\cD_\cP\coloneqq\bigcup_{k,m}\mathcal{A}_{k:m}$.
A configuration $(\dm_1,\dm_2)$ in states $(x_1,x_2)$ is
\emph{fully admissible} if
$(x_i)\in\cD_{S_i}$ and $(\iota(\dm_i))\in\cD_\cP$.
\end{definition}

\begin{theorem}[Operator grammar]\label{thm:grammar}
Let $(\dm_1,\dm_2)$ be fully admissible. The sequence
$L_3\to L_1\to g_1\to L_2\to g_2\to R_1\to R_2\to R_3\to U_1\to U_2\to U_3$
is well-defined and produces a dm3 system at every stage.
Ordering constraints:
(1)~$L_3$ before $g_1$ (separation before expansion);
(2)~$L_1$ before $L_2$ (local before global coherence);
(3)~$g_1$ before $g_2$ (base before recursive expansion);
(4)~$R_1$ before $U_1$ (resonance needed for pushout);
(5)~$R_3$ before $U_3$ (resonance stabilized before synthesis).
Steps involving $L_1$ and $g_2$ are conditional on
Conjectures~\ref{conj:L1} and~\ref{conj:torus}.
The reduced sequence $L_3\to g_1\to R_1\to R_2\to R_3\to U_1\to U_2\to U_3$
holds unconditionally.
\end{theorem}

\begin{proof}
dm3 membership at each stage:
After $L_3$: Proposition~\ref{prop:L3}.
After $L_1$: Proposition~\ref{prop:L1} (conditional).
After $g_1$: Proposition~\ref{prop:g1expand}.
After $L_2$: Proposition~\ref{prop:L2}.
After $g_2$: semi-conjugacy preserved (conditional on Conjecture~\ref{conj:torus}).
After $R_1$: detection only.
After $R_2$, $R_3$: Propositions~\ref{prop:R2},~\ref{prop:R3}.
After $U_1$: Proposition~\ref{prop:pushout}.
After $U_2$: correspondence, no dm3 alteration.
After $U_3$: Proposition~\ref{prop:closure}.
Admissibility maintained since $B_1,B_3$ persist under controlled
operator perturbations (shown case by case above).
\end{proof}

%%=============================================================
\section{Contact geometric physics}\label{sec:contact}
%%=============================================================

\subsection{Why contact geometry}

Seven structural features of dm3 systems force the contact framework:
(1)~hyperbolic limit cycle attractors; (2)~strictly decreasing Lyapunov
functions; (3)~gradient-like operators; (4)~noise-as-fuel stochastic
stability; (5)~embodiment threshold as bifurcation parameter;
(6)~Arnold tongue resonance; (7)~categorical pushouts of dissipative
systems. Features (1)--(2) exclude Hamiltonian mechanics (Liouville's
theorem). Features (4)--(5) require a dissipation parameter evolving
along trajectories. Contact geometry is the unique framework satisfying
all seven simultaneously.

\subsection{The contact manifold}

\begin{definition}[Extended state space and contact form]\label{def:contact-mfld}
Define $M=S\times\R$ with coordinates $(s,z)$, where $z$ is the
\emph{action variable}. Let $\lambda\in\Omega^1(S)$ with $d\lambda=\omega$
a closed non-degenerate 2-form. Set $\alpha=dz-\lambda$.
\end{definition}

\begin{proposition}[Contact structure and Reeb field]\label{prop:contact}
$(M,\alpha)$ is a contact manifold with $R_\alpha=\partial_z$.
\end{proposition}
\begin{proof}
$\alpha\wedge(d\alpha)^n=(dz-\lambda)\wedge(-\omega)^n\ne 0$ since
$\omega^n\ne 0$ and $dz$ is independent of $S$-forms. For Reeb:
$\alpha(\partial_z)=1$ and $\iota_{\partial_z}(-\omega)=0$.
\end{proof}

\subsection{Winding integral and Theorem~A claim (iv)}

\begin{proposition}[Topological invariance of $\oint_\Orb\lambda$]
\label{prop:winding}
\begin{enumerate}[label=\textup{(\roman*)}]
  \item $\oint_\Orb\lambda$ is independent of the choice of $\lambda$
        with $d\lambda=\omega$.
  \item It is invariant under admissible operators fixing $\Orb$ or
        mapping it to a homologous cycle.
  \item It is nonzero when $\Orb$ bounds no chain in $S$.
\end{enumerate}
\end{proposition}
\begin{proof}
(i) $\oint_\Orb df=0$ for any $f$ since $\Orb$ is closed.
(ii) A contactomorphism changes $\lambda$ by an exact form along $\Orb$.
(iii) Stokes' theorem.
\end{proof}

\subsection{Contact Hamiltonian and dynamics}

\begin{definition}[Generative Hamiltonian and equations]\label{def:H}
$H(s,z)=H_{\mathrm{phase}}(s)+H_{\mathrm{diss}}(s,z)$,
where $H_{\mathrm{phase}}=\tfrac12\omega(Y,Y)$ and
$H_{\mathrm{diss}}=-\gamma V(s)e^{-\beta z}$ for $\gamma,\beta>0$.
The full dynamics are $\dot x=X_H(x)+X_{\cR}(x)$, where
$\cR=\tfrac{c}{2}\norm{\dot s-Y(s)}_g^2$ is the Rayleigh function
and $X_{\cR}(s)=-c(\dot s-Y(s))$.
\end{definition}

\begin{proposition}[Recovery of dm3 dynamics]\label{prop:recovery}
$\pi_*X_H(s)=Y(s)-\gamma e^{-\beta z}\nabla V(s)+O(V)$.
On $\Orb$: $\pi_*X_H=Y$.
Along trajectories: $\dot V\le -c\norm{\nabla V}^2\le 0$,
with equality iff $s\in\Orb$.
\end{proposition}

\subsection{Action functional and Euler--Lagrange equations}

\begin{definition}[Generative action]\label{def:action}
$\cS[\gamma]=\int_0^T(\dot z-\lambda(\dot s)-H(s,z))\,dt$
for $\gamma(t)=(s(t),z(t))\colon[0,T]\to M$.
\end{definition}

\begin{theorem}[Contact Euler--Lagrange equations]\label{thm:EL}
Critical points of $\cS$ satisfy
\[
  \dot s=\frac{\partial H}{\partial p}+X_{\cR},\quad
  \dot z=\lambda(\dot s)+H,\quad
  \dot p=-\frac{\partial H}{\partial s}+p\frac{\partial H}{\partial z}.
\]
The contact term $p\,\partial H/\partial z$ is absent from standard
Hamiltonian mechanics.
\end{theorem}

\begin{proposition}[On-orbit action]\label{prop:orbit-action}
$\cS[\gamma|_\Orb]=\oint_\Orb\lambda$ (the topological invariant
of Proposition~\ref{prop:winding}).
\end{proposition}

\subsection{Contact Noether theorem and winding number conservation}

\begin{theorem}[Contact Noether theorem]\label{thm:noether}
If $\{\Phi_\eps\}$ is a one-parameter family of contactomorphisms
with generator $W$ satisfying $\mathcal{L}_W\alpha=f\alpha$ and
$\mathcal{L}_W H=fH$, then $J_W=\alpha(W)$ satisfies
$\dot J_W=f\cdot J_W$ along $X_H$-trajectories.
For strict contactomorphisms ($f=0$): $J_W$ is conserved.
\end{theorem}

\begin{corollary}[Winding number conservation]\label{cor:winding}
Phase rotation on $\Orb$ is a strict contactomorphism. Its conserved
quantity is $\oint_\Orb\lambda$---the topological invariant making
$\Orb$ structurally stable.
\end{corollary}

\subsection{Boundaries as contact submanifolds}

\begin{proposition}[$B_1$ and $B_3$ in contact geometry]\label{prop:contact-bdy}
$\hat B_1=B_1\times\R\subset M$ is a contact hypersurface
($\alpha|_{\hat B_1}$ is a contact form).
$\hat B_3\subset M$ is a Legendrian submanifold
($\alpha|_{\hat B_3}=0$, since $H|_{B_3}=0$).
\end{proposition}

\subsection{Proof of Theorem~C}

\begin{proof}[Proof of Theorem~\ref{thm:C}]
\emph{Step 1 (Floquet coordinates).}
By Floquet theory~\cite{GH}, there exist coordinates $(\rho,\theta)$
near $\Orb$ in $S$ with $\dot\rho=\mu_{\max}\rho+O(\rho^2)$ and
$\dot\theta=\omega+O(\rho)$.

\emph{Step 2 (Contact extension).}
Extend to $M=S\times\R$ with $H$ from Definition~\ref{def:H}.
The contact equations via Theorem~\ref{thm:EL} yield:
$\dot\rho=\mu_{\max}(1-e^{-\beta z})\rho+O(\rho^2)$,
$\dot\theta=\omega+O(\rho)$,
$\dot z=\omega-\abs{\mu_{\max}}\rho^2 e^{-\beta z}+O(\rho^3)$.

\emph{Step 3 (Contactomorphism).}
The coordinate change from the original dm3 coordinates is a
contactomorphism: it preserves $\alpha$ up to a nonvanishing
factor $e^f$, preserving all contact structure.

\emph{Step 4 (Uniqueness).}
The normal form depends only on $(\mu_{\max},\omega,\beta)=\iota(\dm)$.
Any two dm3 systems with the same canonical triple are locally
contact-diffeomorphic.
\end{proof}

%%=============================================================
\section{Structural stability}\label{sec:stability}
%%=============================================================

\subsection{Single-operator and boundary persistence}

\begin{proposition}[Persistence under operators]\label{prop:persist}
For sufficiently small operator parameters:
$g_i(\dm)\in\dmcat$;
$L_i(\dm)\in\dmcat$;
$R_i(\dm_1,\dm_2)\in\dmcat\times\dmcat$;
$U_i(\dm_1,\dm_2)\in\dmcat$.
\end{proposition}
\begin{proof}
Each operator is $C^1$-close to the identity or a dm3 morphism
near $\Orb$; hyperbolicity (open $C^1$ condition) and Lyapunov
descent (continuous in $C^1$) are preserved. See
Propositions~\ref{prop:g1expand}, \ref{prop:L3}, \ref{prop:R3},
\ref{prop:closure}.
\end{proof}

\begin{proposition}[Boundary persistence]\label{prop:bdy-persist}
Under sufficiently small $C^1$ perturbations: $B_1'$, $B_2'$, $B_3'$
exist and vary $C^1$-smoothly.
\end{proposition}
\begin{proof}
$B_1$: stable manifold theorem applied to perturbed $\Lambda'$.
$B_2$: implicit function theorem on $kT_1^*-mT_2^*=0$.
$B_3$: regular level set theorem on $U_{12}^\eps$.
\end{proof}

\subsection{Proof of Theorem~D}

\begin{proof}[Proof of Theorem~\ref{thm:D}]
\emph{Step 1.} Each operator produces a $C^1$-perturbation of $Y$
of size $\delta_i$ (bounded by its parameter; see
Definition~\ref{def:g1} for $g_1$: $O(\delta)$; $L_1,L_3$: $O(\eps)$,
$O(\eta)$; $R_2,R_3$: $O(\norm{W_{k:m}}_{C^1})$).

\emph{Step 2.} The composed perturbation satisfies
$\norm{Y_N-Y}_{C^1}\le\prod_{i=1}^N(1+\delta_i)-1$.
Choosing parameters so this product satisfies
$\prod(1+\delta_i)-1<\varepsilon_0$ keeps the composition within
the $\varepsilon_0$-ball.

\emph{Step 3.} Within the $\varepsilon_0$-ball, Shub's theorem~\cite{Shub}
gives a unique hyperbolic $\Orb'$ near $\Orb$ with
$d_{C^1}(\Orb',\Orb)=O(\varepsilon_0)$. The winding number is preserved
by Proposition~\ref{prop:winding}.

\emph{Step 4.} $T^{*\prime}$, $\mu_{\max}'$, $\tau'$ vary by
$O(\varepsilon_0)$ (period: implicit function theorem; Floquet:
spectral perturbation; threshold: continuity of $\sqrt{c/\kappa}$).

\emph{Step 5.} Proposition~\ref{prop:bdy-persist}.

\emph{Step 6.} $\varepsilon_0<\varepsilon_0$ is satisfied throughout,
so Axioms~\ref{ax:orbit}--\ref{ax:hyperbolic} hold for $\dm'$
by Steps~3--5.
\end{proof}

\begin{remark}[Stability radius]\label{rem:radius}
$\varepsilon_0=\tfrac12\abs{\mu_{\max}}\cdot\sup_\Orb\norm{\Hess V}^{-1}$:
larger $\abs{\mu_{\max}}$ (faster contraction) and flatter $V$ both
enlarge the stability radius. This bound is computable from
$\iota(\dm)$ and the Lyapunov geometry.
\end{remark}

%%=============================================================
\section{Discussion}\label{sec:discussion}
%%=============================================================

\subsection{Summary of main results}

This paper develops a geometric framework for dissipative dynamical
systems with structured limit cycles. The dm3 system, formalized
through eight axioms, is the central object. An operator algebra
acting on the category $\dmcat$ governs expansion ($g$-operators),
coherence ($L$-operators), resonance ($R$-operators), and unification
($U$-operators). The entire structure is embedded in contact geometry
on $M=S\times\R$.

The four main theorems establish: (A)~existence and exponential
stability of the dm3 orbit, with topological invariant
$\oint_\Orb\lambda$ and stochastic stability below $\tau$;
(B)~closure of $\dmcat$ under $\cU$, with the prediction
$\tau_{12}\le\min(\tau_i)$; (C)~the universal contact normal form
with parameters $(\mu_{\max},\omega,\beta)$; and (D)~$C^1$-structural
stability with explicit radius $\varepsilon_0$.

\subsection{Relation to nonequilibrium thermodynamics}

The action variable $z$ plays the role of entropy: it increases
monotonically ($\dot z>0$ on $\Orb$) and the contact form
$\alpha=dz-\lambda$ encodes the first and second laws simultaneously.
The embodiment threshold $\tau$ functions as a noise-induced phase
transition. The result $\tau_{12}\le\min(\tau_i)$ says that
higher-order organization requires lower effective temperature,
consistent with observations in biological and physical oscillator
networks and with stochastic thermodynamic
theory~\cite{Hasminski,PRK}.

\subsection{Relation to coupled oscillator theory}

Period resonance, Arnold tongues, and mode locking are
classical~\cite{PRK,Izhikevich}. The contribution here is embedding
these phenomena in a categorical structure: resonance is the
precondition for the pushout in $\dmcat$. The prediction
$\tau_{12}\le\min(\tau_i)$ is new and testable.

\subsection{Open problems}

\begin{enumerate}[label=(\arabic*),leftmargin=*]
\item \textbf{Global basin expansion} (Conjecture~\ref{conj:global}):
  conditions on $S$ for global (not just local) basin expansion.

\item \textbf{Invariant torus stability} (Conjecture~\ref{conj:torus}):
  KAM theory (quasi-periodic) and Poincar\'e--Melnikov
  (resonant) in the abstract setting. Proved in~\cite{Paper2}
  for the toy model.

\item \textbf{Higher-order resonance and Devil's staircase}:
  measure-theoretic and fractal properties of the full Arnold
  tongue diagram in $\cD_\cP$.

\item \textbf{Infinite-dimensional extensions}: Hilbert/Banach
  manifold generalizations; infinite-dimensional Floquet theory;
  PDE and field-theory dm3 systems.

\item \textbf{Categorical completeness of $\dmcat$}: does $\dmcat$
  have all finite limits? colimits? Is it a topos?
\end{enumerate}

\subsection{Companion paper}

The companion paper~\cite{Paper2} instantiates every definition and
operator on the explicit system
$\dot r=r(1-r^2)+2(r-1)e^{-z}$,
$\dot\theta=1$,
$\dot z=r^2-2(r-1)^2e^{-z}$
on $S=\R^2_{>0}$, verifying all theoretical predictions including
the global attractor, four bifurcations, stationary SDE measure,
and Conjecture~\ref{conj:torus} in the $1$:$2$ resonant case.

%%=============================================================
\appendix
%%=============================================================

\section{The category \texorpdfstring{$\dmcat$}{dm3}: formal development}
\label{app:cat}

[Full categorical development: objects, morphisms, composition law,
associativity, identity, pushouts (proved in
Proposition~\ref{prop:pushout}), coproducts (disjoint union when
$R_1=\emptyset$), and initial/terminal objects if they exist.
Approximately 3--4 pages.]

\section{Supporting lemmas}\label{app:lemmas}

[Proofs deferred from the main text. Approximately 2--3 pages.]

%%=============================================================
\section*{Acknowledgments}
The author acknowledges with gratitude the foundational influence
of the teachings of Paramahamsa Nithyananda and the yogic
scriptural tradition from which this work derives. The mathematical
formalization presented here is understood by the author as a
translation of principles encountered through that tradition.

\begin{thebibliography}{99}
%%=============================================================

\bibitem{Arnold}
V.~I.~Arnold,
\emph{Mathematical Methods of Classical Mechanics},
2nd ed., Springer, New York, 1989.

\bibitem{Bravetti2017}
A.~Bravetti,
\emph{Contact Hamiltonian mechanics},
Ann.\ Phys.\ \textbf{376} (2017), 17--39.

\bibitem{deLeonLainz2019}
M.~de~Le{\'o}n and M.~Lainz~Valc{\'a}zar,
\emph{Contact Hamiltonian systems},
J.\ Math.\ Phys.\ \textbf{60} (2019), 102902.

\bibitem{Gaset2020}
J.~Gaset, X.~Gr{\`a}cia, M.~C.~Mu{\~n}oz-Lecanda, X.~Rivas,
and N.~Rom{\'a}n-Roy,
\emph{New contributions to the Hamiltonian and Lagrangian contact
formalisms for dissipative mechanical systems and optimal control},
J.~Geom.\ Mech.\ \textbf{12} (2020), 543--574.

\bibitem{Geiges}
H.~Geiges,
\emph{An Introduction to Contact Topology},
Cambridge University Press, Cambridge, 2008.

\bibitem{GH}
J.~Guckenheimer and P.~Holmes,
\emph{Nonlinear Oscillations, Dynamical Systems, and Bifurcations
of Vector Fields},
Springer, New York, 1983.

\bibitem{Hasminski}
R.~Z.~Has{'}mi{\u{n}}ski{\u{\i}},
\emph{Stochastic Stability of Differential Equations},
Sijthoff \& Noordhoff, Alphen aan den Rijn, 1980.

\bibitem{HPS}
M.~W.~Hirsch, C.~C.~Pugh, and M.~Shub,
\emph{Invariant Manifolds},
Lecture Notes in Mathematics 583, Springer, Berlin, 1977.

\bibitem{Izhikevich}
E.~M.~Izhikevich,
\emph{Dynamical Systems in Neuroscience},
MIT Press, Cambridge, MA, 2007.

\bibitem{Lee}
J.~M.~Lee,
\emph{Introduction to Smooth Manifolds},
2nd ed., Springer, New York, 2013.

\bibitem{Morrison1986}
P.~J.~Morrison,
\emph{A paradigm for joined Hamiltonian and dissipative systems},
Physica D \textbf{18} (1986), 410--419.

\bibitem{PRK}
A.~Pikovsky, M.~Rosenblum, and J.~Kurths,
\emph{Synchronization: A Universal Concept in Nonlinear Sciences},
Cambridge University Press, Cambridge, 2001.

\bibitem{Shub}
M.~Shub,
\emph{Global Stability of Dynamical Systems},
Springer, New York, 1987.

\bibitem{Paper2}
Pablo Nogueira Grossi,
\emph{The dm3 Operator: Explicit Toy Model and Global Dynamical
Analysis},
companion paper, in preparation.

\end{thebibliography}

\end{document}
